% =====================================================================
%  MSACL DS301 Deep Learning · Lab 2 · Code Hint Sheet
%  Multiple-choice code options for the five fill-in cells.
%  Style matches labs/handouts/lab01_hints.tex.
% =====================================================================
\documentclass[11pt]{article}

\usepackage[letterpaper, margin=0.55in, top=0.55in, bottom=0.4in]{geometry}
\usepackage{amsmath, amssymb}
\usepackage{xcolor}
\usepackage{fancyhdr}
\usepackage{enumitem}

\pagestyle{fancy}
\fancyhf{}
\fancyhead[L]{\small\textbf{MSACL DS301 --- Deep Learning}}
\fancyhead[R]{\small\textbf{Lab 2: Code Hint Sheet}}
\fancyfoot[C]{\thepage}
\renewcommand{\headrulewidth}{0.4pt}

\setlength{\parindent}{0pt}
\setlength{\parskip}{1.5pt}
\newcommand{\opt}[1]{\texttt{#1}}
\setlist[itemize]{leftmargin=2.2em, itemsep=0pt, topsep=0pt, parsep=0pt}

\begin{document}
\small

\begin{center}
{\Large \textbf{Lab 2 --- Code Hint Sheet}}\\[2pt]
{\normalsize A 1D CNN on Real Clinical Spectra}
\end{center}

\vspace{-2pt}
\begin{center}
\fbox{\parbox{0.92\textwidth}{\small
Stuck on a \textbf{YOUR TURN} cell? Each blank you need is below with a few options.
Pick one, type it in, and \textbf{run the cell} --- the built-in \texttt{assert} checks will tell you
whether it is right. Only the five marked cells need editing; run everything else as-is.
}}
\end{center}

\vspace{2pt}\hrule\vspace{4pt}

\textbf{Step 3 --- the training loop.}\; {\small Write these five lines \emph{in order} inside the
batch loop (Lecture 2: forward $\to$ loss $\to$ zero $\to$ backward $\to$ step) --- same recipe
as Lab 1, reused here for spectra.}

\textbf{Blank 1 --- the five-line loop body.}
\begin{itemize}
  \item[(a)] \opt{pred = model(X\_batch)} \quad then \quad \opt{loss = loss\_fn(pred, y\_batch)} \quad then\\
    \opt{optimizer.zero\_grad()} \quad then \quad \opt{loss.backward()} \quad then \quad \opt{optimizer.step()}
  \item[(b)] \opt{pred = loss\_fn(X\_batch)} \quad then \quad \opt{loss = model(pred, y\_batch)} \quad \ldots
  \item[(c)] \opt{loss.backward()} \quad then \quad \opt{pred = model(X\_batch)} \quad then \quad \opt{optimizer.step()} \quad \ldots
\end{itemize}
{\small (a) is the Lecture 2 order; (b) swaps what \texttt{model} and \texttt{loss\_fn} do; (c) backs
up before making a guess --- there is nothing to backprop yet.}

\vspace{3pt}\hrule\vspace{4pt}

\textbf{Step 4 --- the first conv block.}\; {\small One input channel (a raw spectrum), 8 filters,
a 15-bin-wide filter, and padding so the length is unchanged (Lecture 5's code slide).}

\textbf{Blank 2 --- \texttt{self.conv1}.}
\begin{itemize}
  \item[(a)] \opt{nn.Conv1d(in\_channels=1, out\_channels=8, kernel\_size=15, padding=7)}
  \item[(b)] \opt{nn.Conv1d(in\_channels=8, out\_channels=1, kernel\_size=15, padding=7)}
  \item[(c)] \opt{nn.Linear(in\_features=1, out\_features=8)}
\end{itemize}

\vspace{3pt}\hrule\vspace{4pt}

\textbf{Step 4 --- the second conv block.}\; {\small Mirror block 1's pattern, but now reading
the \textbf{8} channels block 1 produced and turning them into \textbf{16}.}

\textbf{Blank 3 --- \texttt{self.conv2}, \texttt{self.relu2}, \texttt{self.pool2}.}
\begin{itemize}
  \item[(a)] \opt{nn.Conv1d(in\_channels=8, out\_channels=16, kernel\_size=15, padding=7)}\\
    \hspace{1.4em}{\small then, same as block 1:}\quad\opt{nn.ReLU()}\quad{\small and}\quad\opt{nn.MaxPool1d(4)}
  \item[(b)] \opt{nn.Conv1d(in\_channels=16, out\_channels=8, kernel\_size=15, padding=7)}
  \item[(c)] \opt{nn.Conv1d(in\_channels=8, out\_channels=16, kernel\_size=1, padding=0)}
\end{itemize}

\vspace{3pt}\hrule\vspace{4pt}

\textbf{Step 4 --- the input size of the first \texttt{Linear} layer.}\; {\small The \texttt{Conv1d}
layers keep the length unchanged (that's the padding); each \texttt{MaxPool1d(4)} divides the length
by 4, and there are \textbf{two} of them; \texttt{Flatten} then stacks all \textbf{16} channels into
one vector. So: $16 \times (6000 \div 4 \div 4)$.}

\textbf{Blank 4 --- \texttt{self.fc1}.}
\begin{itemize}
  \item[(a)] \opt{self.fc1 = nn.Linear(16 * (6000 // 4 // 4), 32)}
  \item[(b)] \opt{self.fc1 = nn.Linear(16 * 6000, 32)}
  \item[(c)] \opt{self.fc1 = nn.Linear(8 * 375, 32)}
\end{itemize}
{\small (a) is right: two \texttt{MaxPool1d(4)}s take 6000 $\to$ 1500 $\to$ 375, and 16 channels give
$16 \times 375 = 6000$ features. (b) forgets the pooling shrank the length; (c) uses the wrong channel
count (8, not 16) --- both make \texttt{fc1} the wrong size and the forward pass will error.}

\vspace{3pt}\hrule\vspace{4pt}

\textbf{Step 5 --- pick a batch size.}\; {\small A middle-ground value for 590 training spectra
(Lecture 4's Goldilocks range).}

\textbf{Blank 5 --- \texttt{CNN\_BATCH\_SIZE}.}
\begin{itemize}
  \item[(a)] \opt{CNN\_BATCH\_SIZE = 32}
  \item[(b)] \opt{CNN\_BATCH\_SIZE = 1}
  \item[(c)] \opt{CNN\_BATCH\_SIZE = 590}
\end{itemize}
{\small (a) and other values in roughly 8--64 all pass the check and train fine --- (b) updates
extremely often but painfully slowly, (c) is the whole training set at once, no mini-batching at all.}

\vspace{3pt}\hrule\vspace{4pt}

\textbf{Step 8 --- commit to a threshold.}\; {\small From your Step 7 exploration, pick a cut that
catches at least \emph{half} the resistant cases (recall $\geq 0.5$). On this trained model recall
stays 0 near \texttt{0.5} and climbs only for low cuts.}

\textbf{Blank 6 --- \texttt{chosen\_threshold}.}
\begin{itemize}
  \item[(a)] \opt{chosen\_threshold = 0.05}
  \item[(b)] \opt{chosen\_threshold = 0.5}
  \item[(c)] \opt{chosen\_threshold = 0.9}
\end{itemize}
{\small (a) reaches recall $\approx 0.67$ --- it clears the \texttt{assert recall\_at\_cut >= 0.5}.
(b) is the default where recall $\approx 0$; (c) is even stricter, still $0$. Any low cut that hits
recall $\geq 0.5$ passes --- there is no single "right" number, it's the clinical trade-off.}

\vspace{4pt}\hrule\vspace{3pt}
{\small\itshape Answers self-check in the notebook: the shape cell confirms \texttt{fc1} accepts the
right number of features and the forward pass returns one logit per spectrum; the batch-size and
threshold cells print their own success messages. Watch the \textbf{resistant recall} number, not just
accuracy --- a model that always says "susceptible" can still score $\sim$94\% accuracy while catching
zero real cases.}

\end{document}
